%=============================================================================
% DATA: REPRODUCTION, AUDIT, CORRECTED DATASET
%=============================================================================
\section{Data}
\label{sec:data}

\subsection{Reproducing the benchmark}
\label{sec:repro}
We ran the reference implementation~\cite{weng2024graph} unmodified---their code,
their array, their protocol---on a pinned software stack. It recovers their
published numbers to within \textbf{7.3\%} across twenty comparisons, with
AAGCN and DCRNN matching to four decimal places. The published results are
real, but three properties of the evaluation account for the reported
advantage. The column labelled ``cross validated'' is computed on a loader
spanning the full series, so it includes training data. RMSE is averaged per
window rather than pooled, which on this target is $\sim$39\% lower. No naive
baseline is reported, and adding persistence to that same column reverses the
ranking: it beats \emph{every} model on MAE and RMSE. An independent
re-implementation that shares no code with either codebase reproduces this.

\subsection{Auditing the benchmark array}
\label{sec:audit}
The benchmark array (\texttt{sri\_lanka\_2013-2022\_shifted.npy}) has no date
axis. We dated every row against the 552 weekly epidemiological reports it was
built from, and dated the climate channels against ERA5 reanalysis and the
authors' own raw satellite files. Table~\ref{tab:audit} lists what we found.

\begin{table}[t]
\centering
\caption{Audit of the benchmark array. Each finding is confirmed against the
source reports or an independent reanalysis.}
\label{tab:audit}
\small
\setlength{\tabcolsep}{3pt}
\begin{tabular}{p{0.33\linewidth}p{0.57\linewidth}}
\toprule
finding & consequence \\
\midrule
Rows 0--47 are 2023; the rest run 2013--2022 & every training split of the
standard protocol contains data from \emph{after} its test period \\
Case column reordered, climate not & cases sit beside weather from a different
week, drifting to 48--55 weeks apart by 2019 \\
``Lagged'' precipitation and canopy channels & match ERA5 best 12--13 weeks
\emph{later}; min-NDVI best 9--18 weeks later: future information \\
49 source reports missing, one week duplicated & 2014 half missing; 2022 stops
in March \\
Week-395 ``19$\times$ spike'' & a formula dragged across one printed table
(15 of 24 cells are each the sum of the two before); the same table's cumulative row gives the
true count \\
4\% of weather values are zero-filled & Jaffna's columns are all zero: it gets no
grid cells in the authors' aggregation \\
\bottomrule
\end{tabular}
\end{table}

Two of these are information leaks. Training on 2023 while testing on
2018--2022 lets a model see the future of its test period. ``Lagged'' climate
channels that carry rainfall and vegetation from three to four months ahead
give it that future's weather. Any score on this array, including the
published ones and our own earlier work, is therefore an optimistic bound.
Earlier drafts of this study called the week-395 spike an administrative
reporting backlog. It is a spreadsheet error, and it can be corrected from the
same published table.

\subsection{The corrected dataset}
\label{sec:corrected}
We rebuilt the case series from every source report, keyed by report volume and
number: 559 weeks from 2013-W26 to 2024-W10. The seven weeks with no report stay
missing and are never interpolated. All 451 weeks the array shares with the
rebuilt series match exactly, so the array's values were correct and only their
order and coverage were wrong. Covariates are aligned row for row: ERA5
temperature, precipitation, humidity and soil moisture for each district;
MODIS NDVI taken \emph{as of} each week, so a composite counts only after its
release; and population using the previous year's official figure.

Lags are applied in the model, not baked into the data: cases up to $t-1$,
climate up to $t-2$, and NDVI and population as of $t$. A test perturbs every
input value at or after each limit and asserts that no model input moves, and
every new input path in this study received its own such test. The
2022--23 serosurvey and serotype-switch dates are used for validation only.
Unless stated otherwise, every number from here on is on this dataset. The
persistence test RMSE on the frozen three origins moves from 44.80 (29.52 with
the week-395 windows excluded) on the array to \textbf{36.02}. The last test
fold now covers the 2023 epidemic.
